\begin{textsample}{NEG Dim 3 – 2009\_c4 – Score 1.00 – 2009\_c4/enade\_2009\_proc\_ger\_1f5954ceea.txt}  \label{ex:f3_neg_160}
Capítulo 4 Percepção sobre a Prova As percepções dos estudantes sobre a prova aplicada no ENADE/2009 na área de TECNOLOGIA EM PROCESSOS GERENCIAIS foram mensuradas por meio de nove itens que avaliaram desde o grau de dificuldade da prova até o tempo gasto para concluí -la.
As questões foram analisadas separando -se concluintes e ingressantes, e as percepções sobre a prova foram relacionadas com o desempenho dos estudantes e com a Região Geográfica de localização do curso6.
O desempenho dos estudantes dentro de cada Área/Subárea foi classificado em quartos ( 1/4 ) para os concluintes e para os ingressantes em separado ( ver Figura 1 para o esquema de representação gráfica destas grandezas ).
As notas finais foram colocadas em ordem ascendente e no primeiro quarto foram colocadas as 25 \% notas mais baixas, limitadas superiormente pelo percentil 25 que vem a ser o primeiro quartil.
No segundo quarto, as 25 \% seguintes, limitadas superiormente pelo percentil 50 que vem a ser a mediana.
No terceiro quarto, as 25 \% notas após as encontradas no segundo quarto, ou seja, as primeiras 25 \% acima da mediana, limitadas superiormente pelo percentil 75 que vem a ser o terceiro quartil.
No último quarto foram alocadas as 25 \% maiores notas.
Potencialmente, poderia existir uma relação entre a nota do estudante e a sua percepção sobre a prova.
Calculando as médias das respostas segundo os quartos de notas, seria possível reconhecer a existência ( ou não ) de algum gradiente nas respostas, por exemplo, uma melhor nota correspondendo a uma percepção menos negativa.
Este gradiente seria traduzido, por exemplo no gráfico 4.1a, pela diminuição das áreas laranja e vermelha ( Difícil e Muito difícil ) com o aumento do desempenho ( indo do primeiro ao último quarto ).
A medida tradicional de “ grau de discriminação ” comparou os resultados dos estudantes do primeiro quarto ( 25 \% piores estudantes ) aos do último quarto ( 25 \% melhores ).
Esta comparação, no gráfico disponibilizado, foi idêntica à comparação da primeira barra ( primeiro quarto ) com a última barra ( último quarto ) de cada grupo ingressante ou concluinte nos gráficos “ segundo quartos de nota final ” ( gráficos com números ímpares neste capítulo ).
Por outro lado, a disponibilização das informações segundo os quartos de nota permitiu verificar o comportamento de dois outros grupos intermediários que na análise somente dos extremos tem que ser suposto monotônico. 6 Curso, neste relatório, refere -se à unidade de análise para o Conceito ENADE e é caracterizado pela combinação de IES, Área/Subárea ( quando pertinente ) e município de habilitação. primeira barra ( primeiro quarto ) com a última barra ( último quarto ) de cada grupo ingressante ou concluinte nos gráficos “ segundo quartos de nota final ” ( gráficos com números ímpares neste capítulo ).
Por outro lado, a disponibilização das informações segundo os quartos de nota permitiu verificar o comportamento de dois outros grupos intermediários que na análise somente dos extremos tem que ser suposto monotônico.
A seguir, serão apresentados os principais resultados relativos aos nove itens avaliados. 4.1 Grau de dificuldade da prova 4.1.1 Formação Geral Os Gráficos 4.1a, 4.1b e 4.1c apresentam a distribuição das respostas à questão “ Qual o grau de dificuldade desta prova na parte de Formação Geral? ”, segundo quartos de nota final e grupo de estudantes concluintes e ingressantes que participaram do ENADE/2009 na área de TECNOLOGIA EM PROCESSOS GERENCIAIS.
No total de concluintes, 30,2 \% declararam que a prova na parte da Formação Geral foi Difícil/Muito difícil, contra um valor próximo para os ingressantes, 32,3 \%.
Tanto entre os concluintes, como entre os ingressantes, a proporção de indivíduos que classificaram esta parte como Difícil/Muito difícil diminuiu com o desempenho do estudante.
Por outro lado, a proporção que declarou a parte Fácil/Muito fácil foi menor entre os concluintes ( 7,8 \% ) do que entre os ingressantes ( 8,1 \% ).
Entre os concluintes, esta proporção aumentou com o desempenho e 62,0 \% classificaram esta parte com um nível Médio de dificuldade.
A proporção entre os ingressantes foi menor, 59,5 \%.
Entre os concluintes, esta proporção aumentou com o desempenho até o terceiro quarto, com uma pequena queda no grupo de melhor desempenho, indo de 58,6 \% no primeiro quarto de notas a 63,2 \% no último quarto.
Entre os ingressantes esta proporção aumentou com o desempenho, indo de 55,8 \% no primeiro quarto a 62,1 \% no último.
Os Gráficos 4.2a e 4.2b apresentam a distribuição das respostas à questão “ Qual o grau de dificuldade desta prova na parte de Formação Geral? ”, segundo a Região onde o curso está situado para os estudantes concluintes e ingressantes que participaram do ENADE/2009 na área de TECNOLOGIA EM PROCESSOS GERENCIAIS.
Entre os concluintes, a maior proporção de estudantes que considerou a parte de Formação Geral da prova como de dificuldade média entre todas as Regiões brasileiras, foram os da Região Nordeste com 68,3 \% e entre os ingressantes também com 72,0 \%.
A Região com as menores proporções desta resposta foram respectivamente para os concluintes e ingressantes foi a Sul ( 58,2 \% e 55,8 \% ). 4.1.2 Componente Específico Os Gráficos 4.3a, 4.3b e 4.3c apresentam a distribuição das respostas à questão “ Qual o grau de dificuldade desta prova na parte do Componente Específico? ”, segundo quartos de nota final e grupo de estudantes concluintes e ingressantes que participaram do ENADE/2009 na área de TECNOLOGIA EM PROCESSOS GERENCIAIS.
Em geral, para as desagregações consideradas, os estudantes classificaram o Componente Específico como mais difícil que a Formação Geral.
No total de concluintes, 30,6 \% declararam que a prova na parte do Componente Específico foi Difícil/Muito difícil, sendo um valor muito maior para os ingressantes, 34,8 \%.
Entre os concluintes, a proporção de indivíduos que classificaram esta parte como Difícil/Muito difícil entre os concluintes é declinante com o desempenho partindo de 32,5 \% no primeiro quarto para 28,8 \% no último quarto.
Entre os ingressantes esta proporção também decresceu com o desempenho.
A proporção que declarou a parte Fácil/Muito fácil foi maior entre os concluintes ( 7,5 \% ) do que entre os ingressantes ( 6,5 \% ).
Entre os concluintes a proporção dos que consideraram esta parte como fácil se mantém com valores muito próximos a média.
O mesmo comportamento se verificou entre os ingressantes.
Entre os concluintes, 61,9 \% classificaram esta parte com um nível Médio de dificuldade.
A proporção entre os ingressantes foi ligeiramente menor, 58,8 \%.
Entre os concluintes e ingressantes, esta proporção aumentou com o desempenho.
Os Gráficos 4.4a e 4.4b apresentam a distribuição das respostas à questão “ Qual o grau de dificuldade desta prova na parte do Componente Específico? ”, segundo a Região onde o curso está situado para os estudantes concluintes e ingressantes que participaram do ENADE/2009 na área de TECNOLOGIA EM PROCESSOS GERENCIAIS.
Entre os concluintes, a maior proporção de estudantes que considerou a parte do Componente Específico da prova como de dificuldade média entre todas as Regiões brasileiras, foi na Região Norte com 75,0 \%, já entre os ingressantes, foram os estudantes da Região Nordeste com 67,5 \%.
A Região com as menores proporções desta resposta foram para os concluintes e ingressantes, a Sul ( 58,5 \% e 55,1 \% respectivamente ). 4.2 Extensão da prova em relação ao tempo total Os Gráficos 4.5a, 4.5b e 4.5c apresentam a distribuição das respostas à questão sobre a extensão da prova em relação ao tempo total, segundo quartos de nota final e grupo de estudantes concluintes e ingressantes que participaram do ENADE/2009 na área de TECNOLOGIA EM PROCESSOS GERENCIAIS.
Tanto para concluintes quanto para ingressantes, a maioria considerou a extensão da prova como adequada em relação ao tempo destinado à sua resolução, respectivamente 61,6 \% e 60,5 \%.
Entre os concluintes esta proporção cresceu com o desempenho.
Entre os ingressantes esta proporção também foi crescente.
A proporção dos que consideraram a extensão da prova muito longa foi ligeiramente menor entre os concluintes ( 12,9 \% ) do que entre os ingressantes ( 13,7 \% ).
A proporção dos que consideraram a extensão da prova curta/muito curta vis-à-vis ao tempo total foi menor entre os concluintes ( 2,3 \% ) do que entre os ingressantes ( 2,6 \% ).
A maior proporção de ingressantes declarando que a prova foi curta/muito curta foi entre os estudantes do segundo quarto, 2,7 \%.
A proporção dos que consideraram a extensão da prova muito curta foi, entre os concluintes, 0,3 \% e entre os ingressantes, 0,5 \%.
A proporção dos que consideraram a extensão da prova longa foi menor entre os concluintes ( 23,2 \% ) do que entre os ingressantes ( 23,3 \% ).
Os Gráficos 4.6a e 4.6b apresentam a distribuição das respostas à questão sobre a extensão da prova em relação ao tempo total, segundo a Região onde o curso está situado para os estudantes concluintes e ingressantes que participaram do ENADE/2009 na área de TECNOLOGIA EM PROCESSOS GERENCIAIS.
Em relação a percepção sobre o tempo da prova ser adequado, somente os concluintes da Região Nordeste ( 54,0 \% ) e Norte ( 59,6 \% ) se situaram abaixo da média ( 61,6 \% ) para esta resposta, enquanto entre os ingressantes, os da Região Sul ( 61,6 \% ) e Centro-Oeste ( 63,4 \% ) ficaram acima da média ( 60,5 \% ). 4.3 Compreensão dos Enunciados das Questões 4.3.1 Formação Geral Os Gráficos 4.7a, 4.7b e 4.7c apresentam a distribuição das respostas à questão “ Os enunciados das questões da prova na parte de Formação Geral estavam claros e objetivos? ”, segundo quartos de nota final e grupo de estudantes concluintes e ingressantes que participaram do ENADE/2009 na área de TECNOLOGIA EM PROCESSOS GERENCIAIS.
Com relação aos enunciados das questões da parte de Formação Geral, as opiniões foram positivas : a grande maioria dos concluintes ( 73,9 \% ) e ingressantes ( 72,9 \% ) consideraram claros e objetivos os enunciados de todas ou da maioria das questões.
Entre os concluintes, esta proporção cresce com o desempenho : 65,7 \% para o grupo com o pior desempenho e 80,4 \% para o de melhor.
Entre os ingressantes, esta proporção apresentou o mesmo comportamento.
Em linhas gerais, pode -se dizer que a combinação de “ poucos ” e “ não, nenhum ” decresceu com o desempenho tanto para concluintes como para ingressantes.
Os Gráficos 4.8a e 4.8b apresentam a distribuição das respostas à questão “ Os enunciados das questões da prova na parte de Formação Geral, estavam claros e objetivos? ”, segundo a Região onde o curso está situado para os estudantes concluintes e ingressantes que participaram do ENADE/2009 na área de TECNOLOGIA EM PROCESSOS GERENCIAIS.
Em termos regionais, as proporções de ingressantes e de concluintes que consideraram claros e objetivos os enunciados de todas ou da maioria das questões foram respectivamente, 72,9 \% e 73,9 \%.
Entre os concluintes, os valores foram de 73,1 \% no Sul a 78,9 \% no Norte.
Já entre os ingressantes estas proporções variaram de 71,7 \% no Sul a 79,7 \% no Norte.
Com relação às respostas “ poucos ” e “ não, nenhum ”, as maiores proporções ocorreram para concluintes, na Região Sul ( 12,2 \% ) e para ingressantes, na Centro-Oeste ( 14,1 \% ). 4.3.2 Componente Específico Os Gráficos 4.9a, 4.9b e 4.9c apresentam a distribuição das respostas à questão “ Os enunciados das questões da prova na parte de formação específica estavam claros e objetivos? ”, segundo quartos de nota final e grupo de estudantes concluintes e ingressantes que participaram do ENADE/2009 na área de TECNOLOGIA EM PROCESSOS GERENCIAIS.
Com relação aos enunciados das questões da parte do Componente Específico, as opiniões foram também positivas : a grande maioria dos concluintes ( 75,8 \% ) e ingressantes ( 73,4 \% ) consideraram claros e objetivos os enunciados de todas ou da maioria das questões.
A proporção de concluintes declarando os enunciados de todas ou da maioria das questões estavam claros e objetivos foi crescente com o desempenho indo de 68,6 \% para o primeiro quarto para 81,1 \% para os com melhor desempenho.
Entre os ingressantes, esta proporção também cresce, com o desempenho : indo de 67,2 \% para os com pior desempenho a 79,4 \% para os de melhor desempenho.
O número de estudantes que declaram que “ não, nenhum ” dos enunciados das questões da parte do Componente Específico era claro e objetivo, foram semelhantes à da Formação Geral ( 0,8 \% entre os concluintes e 1,0 \% entre os ingressantes ) : 0,8 entre os concluintes e ingressantes.
Os que responderam “ poucos ”, são 8,6 \% entre os concluintes e 9,6 \% entre os ingressantes.
A proporção de estudantes que respondeu que cerca de metade dos enunciados da parte do Componente Específico era claro e objetivo, decresceu com o desempenho, tanto para concluintes quanto para ingressantes.
Os Gráficos 4.10a e 4.10b apresentam a distribuição das respostas à questão “ Os enunciados das questões da prova na parte de formação específica estavam claros e objetivos? ”, segundo a Região onde o curso está situado para os estudantes concluintes e ingressantes que participaram do ENADE/2009 na área de TECNOLOGIA EM PROCESSOS GERENCIAIS.
Em termos regionais, as proporções de ingressantes e de concluintes que consideraram claros e objetivos os enunciados de todas ou da maioria das questões variaram mais entre os ingressantes ( no máximo 8 pontos percentuais ) do que entre os concluintes ( no máximo 5 pontos percentuais ).
Entre os concluintes, os valores foram de 75,4 \% no Sul a 79,8 \% no Norte.
Já entre os ingressantes estas proporções variaram de 71,7 \% na Região Sul a 80,0 \% no Norte.
Com relação às respostas “ poucos ” e “ não, nenhum ” para concluintes a maior proporção ocorreu na Região Centro-Oeste com 11,7 \%.
Para os ingressantes a maior proporção ocorreu na Região Norte com 11,7 \%.
Tanto para concluintes como para ingressantes a opção que apresentou o maior percentual de escolha foi “ sim, a maioria ”, com 57,3 \% e 55,4 \% respectivamente. 4.4 Suficiência das informações/instruções fornecidas Os Gráficos 4.11a, 4.11b e 4.11c apresentam a distribuição das respostas à questão “ As informações/instruções fornecidas para a resolução das questões foram suficientes para resolvê -las? ”, segundo quartos de nota final e grupo de estudantes concluintes e ingressantes que participaram do ENADE/2009 na área de TECNOLOGIA EM PROCESSOS GERENCIAIS.
Quando perguntados sobre as informações/instruções fornecidas para a resolução das questões, quase 82,4 \% dos ingressantes e 84,7 \% dos concluintes de todo o Brasil responderam que estas eram ou excessivas ou suficientes em todas ou na maioria das questões, demonstrando uma avaliação positiva da prova.
Levando -se em conta o desempenho dos participantes, percebem -se diferenças entre as opiniões de estudantes, ingressantes e concluintes.
A proporção de estudantes com a opinião supracitada é crescente com o desempenho, tanto para concluintes como para ingressantes.
Entre os concluintes, a proporção vai de 75,9 \% para os de pior desempenho a 91,3 \% para os de melhor desempenho.
Entre os ingressantes, as proporções equivalentes são, respectivamente, 75,8 \% a 89,3 \%.
O complemento desta opinião, ou seja, os estudantes que responderam somente em algumas ou em nenhuma delas, apresentou uma proporção decrescente com o desempenho tanto para os concluintes como para os ingressantes.
A resposta não, em nenhuma delas, entre os concluintes e ingressantes, apareceu em todos os grupos de desempenho.
Os Gráficos 4.12a e 4.12b apresentam a distribuição das respostas à questão “ As informações/instruções fornecidas para a resolução das questões foram suficientes para resolvê -las? ”, segundo Região onde o curso está situado para os estudantes concluintes e ingressantes que participaram do ENADE/2009 na área de TECNOLOGIA EM PROCESSOS GERENCIAIS.
A Região Nordeste foi a que apresentou o percentual mais baixo, com 82,9 \% entre os concluintes emitindo tal opinião.
E a Região Norte foi a que apresentou o percentual mais baixo, com 79,9 \% dos ingressantes emitindo tal opinião. 4.5 Dificuldade encontrada ao responder à prova Os Gráficos 4.13a, 4.13b e 4.13c apresentam a distribuição das respostas à questão “ Você se deparou com alguma dificuldade ao responder a prova? ”, segundo quartos de nota final e grupo de estudantes concluintes e ingressantes que participaram do ENADE/2009 na área de TECNOLOGIA EM PROCESSOS GERENCIAIS.
Perguntados sobre as dificuldades com as quais se depararam ao responder à prova, 45,3 \% dos ingressantes e 48,8 \% dos concluintes apontaram como principal dificuldade à “ Forma diferente de abordagem do conteúdo ”.
Ingressantes e concluintes apresentam opiniões muito semelhantes no tocante à ordem das opções escolhidas pelo contingente.
Para a quase totalidade das áreas, ingressantes apresentaram uma proporção de estudantes bem maior do que entre os concluintes declarando como principal dificuldade o desconhecimento da matéria.
O desconhehecimento do conteúdo foi apontado somente por 7,4 \% dos concluintes, sendo esta variação entre os quartos de nota final maior entre os concluintes.
A falta de motivação foi considerada como maior dificuldade por 16,3 \% dos ingressantes, e para 16,6 \% dos concluintes.
Em termos nacionais, a proporção de concluintes e ingressantes que declararam não terem tido qualquer dificuldade para responder à prova, respectivamente 21,2 \% e 17,6 \%.
Os Gráficos 4.14a e 4.14b apresentam a distribuição das respostas à questão “ Você se deparou com alguma dificuldade ao responder a prova? ”, segundo a Região onde o curso está situado para os estudantes concluintes e ingressantes que participaram do ENADE/2009 na área de TECNOLOGIA EM PROCESSOS GERENCIAIS.
Em relação às Regiões brasileiras, o percentual de ingressantes que apontaram o desconhecimento do conteúdo como principal dificuldade ao responder à prova variou de 8,5 \% na Região Norte ( o menor valor ), a 18,6 \% na Centro-Oeste ( o maior valor ).
Entre os concluintes, tal percentual variou de 4,6 \% na Região Nordeste a 10,8 \%, na Norte.
A forma diferente de abordagem de conteúdo foi citada como maior dificuldade por 36,3 \% dos concluintes na Região Norte chegando a 51,1 \% na Região Sul.
Já entre os ingressantes, os valores foram de 40,1 \% na Região Nordeste a 47,5 \% na Região Sul.
O percentual de concluintes que citaram a falta de motivação como principal dificuldade variou de 14,7 \% na Região Centro-Oeste a 19,9 \% na Região Nordeste.
Entre os ingressantes, as estatísticas correspondentes foram de 15,3 \% na Região Norte a 21,6 \%, na Região Nordeste.
Os que declararam que não tiveram qualquer tipo de dificuldade variaram, entre os concluintes, de 19,8 \% na Região Sul a 26,8 \% na Região Nordeste.
Entre os ingressantes, este percentual variou de 16,0 \% na Região Sul a 23,7 \% na Região Norte. 4.6 Percepção sobre os conteúdos das questões objetivas da prova Os Gráficos 4.15a, 4.15b e 4.15c apresentam a distribuição das respostas à questão “ Considerando apenas as questões objetivas da prova, você percebeu que :... ” segundo quartos de nota final e grupo de estudantes concluintes e ingressantes que participaram do ENADE/2009 na área de TECNOLOGIA EM PROCESSOS GERENCIAIS.
As respostas possíveis para esta questão estavam relacionadas à cobertura do conteúdo das questões objetivas durante o curso.
A proporção de ingressantes ( 32,2 \% ) que declarou que não estudou ainda a maioria desses conteúdos foi bem maior do que entre os concluintes ( 5,6 \% ).
Esta proporção foi decrescente tanto para os concluintes como para os ingressantes como função do desempenho.
Por outro lado, a proporção de estudantes que declarou que estudou e aprendeu todos ou a maioria desses conteúdos, entre os ingressantes foi de 44,3 \% e entre os concluintes 69,6 \%.
Esta proporção foi crescente como função do desempenho, tanto para os concluintes como para os ingressantes.
A proporção de estudantes que declarou ter estudado a maioria desses conteúdos, mas não os ter aprendido foi bem maior entre os concluintes ( 12,7 \% ) do que os ingressantes ( 7,6 \% ).
Os Gráficos 4.16a e 4.16b apresentam a distribuição das respostas à questão “ Considerando apenas as questões objetivas da prova, você percebeu que :... ”, para os estudantes concluintes e ingressantes que participaram do ENADE/2009 na área de TECNOLOGIA EM PROCESSOS GERENCIAIS, segundo a Região onde o curso está situado.
Em relação às Regiões brasileiras, nenhum concluinte afirmou não ter estudado ainda a maioria dos conteúdos, os percentuais variaram de 4,6 \% na Região Nordeste a 6,7 \%, na Norte.
Entre os ingressantes, os percentuais variaram de 20,3 \% na Região Norte a 35,3 \% na Sul.
Em todas as Regiões, a maioria dos concluintes afirmou ter estudado e aprendido todos ou muitos dos conteúdos, com percentuais que variaram de 66,8 \% na Região Sul a 78,6 \% na Centro-Oeste.
Para os ingressantes, tal percentual variou de 40,3 \%, na Região Sul a 64,4 \%, na Norte. 4.7 Tempo gasto para concluir a prova Os Gráficos 4.17a, 4.17b e 4.17c apresentam a distribuição das respostas à questão “ Qual foi o tempo gasto para concluir a prova? ”, segundo quartos de nota final e grupo de estudantes concluintes e ingressantes que participaram do ENADE/2009 na área de TECNOLOGIA EM PROCESSOS GERENCIAIS.
Quando perguntados quanto ao tempo de conclusão da prova, a maioria respondeu entre, duas e três horas tanto entre os concluintes ( 40,4 \% ) quanto entre os ingressantes ( 40,1 \% ).
A segunda maior incidência de respostas foi entre uma e duas horas para ambos os grupos : concluintes com 33,6 \% e ingressantes com 37,4 \%.
A proporção daqueles que declararam ter usado as quatro horas, mas não terem conseguido terminar foi bem baixa : 1,3 \% entre os ingressantes e entre os concluintes.
Entre os concluintes esta resposta foi maior entre os estudantes do primeiro quarto, aquele com pior desempenho.
Os Gráficos 4.18a e 4.18b apresentam a distribuição das respostas à questão “ Qual foi o tempo gasto para concluir a prova? ” segundo a Região onde o curso está situado para os estudantes concluintes e ingressantes que participaram do ENADE/2009 na área de TECNOLOGIA EM PROCESSOS GERENCIAIS.
A proporção de concluintes que gastaram entre uma e duas horas foi maior na Região Sul ( 34,4 \% ) e menor na Região Norte ( 22,1 \% ).
Para ingressantes, o maior percentual para este tempo foi maior na Região Nordeste ( 38,8 \% ) e a menor na Região Norte ( 21,7 \% ).

% matched lemmas: 
\end{textsample}
